% Generated draft: requires mathematical and source review.
\documentclass[11pt]{article}
\usepackage[margin=25mm]{geometry}
\usepackage{fontspec}
\setmainfont{Latin Modern Roman}
\usepackage{amsmath,amssymb,mathtools,unicode-math}
\setmathfont{Latin Modern Math}
\usepackage{enumitem,needspace,xurl,hyperref}
\setlength{\parindent}{0pt}
\setlength{\parskip}{5pt}
\setlength{\emergencystretch}{3em}
\newcommand{\N}{\mathbb{N}}
\newcommand{\Z}{\mathbb{Z}}
\newcommand{\Q}{\mathbb{Q}}
\newcommand{\R}{\mathbb{R}}
\newcommand{\C}{\mathbb{C}}
\newcommand{\F}{\mathbb{F}}
\newcommand{\A}{\mathbb{A}}
\newcommand{\PP}{\mathbb{P}}
\newcommand{\E}{\mathbb{E}}
\newcommand{\eps}{\varepsilon}
\newcommand{\dd}{\,\mathrm{d}}
\begin{document}
\section*{UnsolvedMath ID: 30005428}
\textbf{Source problem code:} catalogue entry
\subsection{Definitions and mathematical statement}
Let \(A\) be a finite-dimensional algebra over a field \(k\). Write \(\operatorname{proj}A\) for the category of finitely generated projective \(A\)-modules and \(K^b(\operatorname{proj}A)\) for its bounded homotopy category. Its Grothendieck group has realification
\[
K_0(\operatorname{proj}A)_{\R}:=K_0(\operatorname{proj}A)\otimes_{\Z}\R\cong\R^{|A|},
\]
where \(|A|\) denotes the number of isomorphism classes of indecomposable finitely generated projective \(A\)-modules.

A complex \(T\in K^b(\operatorname{proj}A)\) is presilting if
\[
\operatorname{Hom}_{K^b(\operatorname{proj}A)}(T,T[\ell])=0
\qquad\text{for every integer }\ell>0.
\]
It is silting if, additionally, the smallest thick subcategory containing \(T\) is all of \(K^b(\operatorname{proj}A)\). Here a thick subcategory is a full triangulated subcategory closed under direct summands. A complex is two-term if its terms vanish outside degrees \(-1\) and \(0\). For such a complex, its Grothendieck class is \([T]=[T^0]-[T^{-1}]\).

For a two-term presilting complex with indecomposable direct summands \(T=T_1\oplus\cdots\oplus T_r\), put
\[
C(T):=\left\{\sum_{i=1}^{r}a_i[T_i]:a_i\in\R_{\geq0}\right\},
\qquad
\Sigma(A):=\{C(T):T\text{ is a two-term presilting complex of }A\}.
\]
The zero complex contributes the zero cone. The supplied extract calls \(\Sigma(A)\) the \(g\)-fan of \(A\). A fan is complete when the union of its cones is its entire ambient vector space.

\textbf{Requested characterization.} Characterize those complete sign-coherent fans \(\mathcal F\) in \(\R^d\) that are realizable as \(\Sigma(A)\) for finite-dimensional algebras \(A\).

A possible real-linear interpretation of realizability is the existence of a field \(k\), a finite-dimensional \(k\)-algebra \(A\) with \(|A|=d\), and a real-linear isomorphism
\[
\Phi:K_0(\operatorname{proj}A)_{\R}\longrightarrow\R^d
\quad\text{such that}\quad
\mathcal F=\{\Phi(C):C\in\Sigma(A)\}.
\]
This interpretation is conditional: the supplied record does not specify whether the field is fixed, whether \(\Phi\) must preserve an integral lattice or a distinguished cone, or what additional structure must be preserved.

\textbf{Unresolved qualifications.} The extract asserts that \((\Sigma(A),C(A))\) is a sign-coherent fan, but supplies no definition of sign coherence for a fan with a distinguished cone. Consequently, the input class in the requested characterization cannot be defined precisely from the supplied material alone. No definition of sign coherence or convention for equivalence of fans is imposed here.

The extract also reports that \(\Sigma(A)\) is complete if and only if \(A\) is \(\tau\)-tilting finite, meaning that there are only finitely many isomorphism classes of basic two-term silting complexes. Here basic means that the indecomposable direct summands are pairwise nonisomorphic. This equivalence is supplied background, not a solution of the characterization problem. The extracted passage itself does not contain the characterization question; that target comes from the registry's statement and clean statement.
\subsection{Short English statement}
Determine which complete sign-coherent fans arise from the cones spanned by the Grothendieck classes of indecomposable summands of two-term presilting complexes over finite-dimensional algebras. The supplied material leaves sign coherence, the allowed identifications of fans, and the choice of base field insufficiently specified. This is a statement-only draft requiring human mathematical review.
\subsection{Sources}
\begin{itemize}
\item Local UnsolvedMath registry, permanent ID \(30005428\), title: \emph{Realizability of Complete Sign-Coherent Fans}; exact source code: \texttt{catalogue entry}. The supplied record gives proposed year \(2023\), names no proposer, labels the problem open, and labels its statement \texttt{corrected\_verified}. Its statement-verification field claims checking against the official report; this is supplied provenance, not an independent verification here.
\item \emph{Representation Theory of Quivers and Finite-Dimensional Algebras} (2023), Oberwolfach Reports, pp.\ 397--486, DOI \texttt{10.4171/owr/2023/7}: \url{https://doi.org/10.4171/owr/2023/7}. The supplied extract carries the page marker \(473\). This is also the supplied literature-sources link, labelled ``original OWR report.''
\item The extract attributes its Definition 2 to reference labels [H, DIJ] and Proposition 3 to [A, DIJ]. Bibliographic expansions of these labels are not supplied, so no authors or further references are inferred.
\item Supplied literature assessment, checked \texttt{2026-08-22}: the registry labels the problem open and says that the extracted passage announces a problem without including its statement. It also reports unresolved surrounding questions on deformations and semigroup operations. This is supplied dated curation, not independently verified; the supplied background explicitly requests independent verification of time-sensitive status and cited claims.
\end{itemize}
\end{document}
